\documentclass{article}
\usepackage{lmodern}
\usepackage{amsmath}
\usepackage{listings}
\usepackage{fullpage}
\usepackage{parskip}
\usepackage{xcolor}
\lstset{
	basicstyle=\ttfamily,
	columns=fullflexible,
	frame=single,
	breaklines=true,
	postbreak=\mbox{\textcolor{red}{$\hookrightarrow$}\space},
}
\begin{document}
\title{Experiment `p\_4\_6\_find\_minimum\_seq' Results}
\author{\today}
\date{}
\maketitle
\textbf{Experiment outcome: }SUCCESS
\\\textbf{Bad responses: }0
\\\textbf{Responses containing}~\texttt{assume}~\textbf{: }0
\\\textbf{Responses containing}~\texttt{expect}~\textbf{: }0
\\\textbf{Resolution attempts: }1
\\\textbf{Hard fails (resolution): }0
\\\textbf{Soft fails (resolution): }0
\\\textbf{Verification attempts: }1
\section*{Problem Specification}
\textbf{Problem name: }p\_4\_6\_find\_minimum\_seq
\\\textbf{Natural language statement: }Write a method that finds the minimum value from a list of inputs into a method.
\\\textbf{Method signature: }p\_4\_6\_find\_minimum\_seq(inputs: seq<int>) returns (min: int)
\subsection*{Ensures}
\begin{itemize}
\item \texttt{forall i :: 0 <= i < |inputs| ==> min <= inputs[i]}
\item \texttt{exists i :: 0 <= i < |inputs| \&\& min == inputs[i]}
\end{itemize}
\subsection*{Requires}
\begin{itemize}
\item \texttt{0 < |inputs|}
\end{itemize}
\clearpage
\section*{GenAI interactions}
Below you will find all interactions between the `user' (program) and the `assistant' (OpenAI).
\subsection*{Program $\rightarrow$ GenAI}
\begin{lstlisting}
You are given the following task to perform in Dafny:

Write a method that finds the minimum value from a list of inputs into a method.

The signature should be:

p_4_6_find_minimum_seq(inputs: seq<int>) returns (min: int)

The method should respect the following contract:

requires 0 < |inputs|, ensures forall i :: 0 <= i < |inputs| ==> min <= inputs[i], ensures exists i :: 0 <= i < |inputs| && min == inputs[i]

Produce and show only the Dafny body of this method, including the curly braces that surround it. Do not show the signature nor contract. You must not use 'assume' nor 'expect' anywhere in your code.
\end{lstlisting}
\subsection*{GenAI $\rightarrow$ Program}
\textbf{System fingerprint: }null
\\\textbf{ID: }chatcmpl-DmfdHbRhGD2K0OC2aTN4E88Y02saQ
\\\textbf{Model: }gpt-5.5-2026-04-23
\\\textbf{Created at: }1780492151
\\\textbf{Finish reason: }stop
\\\textbf{Usage: }338 tokens in, and 693 tokens out
\begin{lstlisting}
{
  min := inputs[0];
  var i := 1;
  var k := 0;

  while i < |inputs|
    invariant 1 <= i <= |inputs|
    invariant 0 <= k < i
    invariant min == inputs[k]
    invariant forall j :: 0 <= j < i ==> min <= inputs[j]
  {
    if inputs[i] < min {
      min := inputs[i];
      k := i;
    }
    i := i + 1;
  }

  assert 0 <= k < |inputs| && min == inputs[k];
  assert exists j :: 0 <= j < |inputs| && min == inputs[j];
}
\end{lstlisting}
\section*{Final Program}
\begin{lstlisting}
method p_4_6_find_minimum_seq(inputs: seq<int>) returns (min: int)
	requires 0 < |inputs|
	ensures forall i :: 0 <= i < |inputs| ==> min <= inputs[i]
	ensures exists i :: 0 <= i < |inputs| && min == inputs[i]
{
  min := inputs[0];
  var i := 1;
  var k := 0;

  while i < |inputs|
    invariant 1 <= i <= |inputs|
    invariant 0 <= k < i
    invariant min == inputs[k]
    invariant forall j :: 0 <= j < i ==> min <= inputs[j]
  {
    if inputs[i] < min {
      min := inputs[i];
      k := i;
    }
    i := i + 1;
  }

  assert 0 <= k < |inputs| && min == inputs[k];
  assert exists j :: 0 <= j < |inputs| && min == inputs[j];
}
\end{lstlisting}
\section*{Total Token Usage}
\textbf{Input tokens: }338
\\\textbf{Output tokens: }693
\\\textbf{Reasoning tokens: }512
\\\textbf{Sum of `total tokens': }1031
\section*{Experiment Timings}
\textbf{Overall Experiment} started at 1780492150588, ended at 1780492171091, lasting 20503ms (20.50 seconds)
\\\textbf{Iteration \#1} started at 1780492150588, ended at 1780492171091, lasting 20503ms (20.50 seconds)
\end{document}
